\documentclass[11pt,a4paper]{article}
% Common preamble for the hint sheets (\input from each file)
\usepackage[utf8]{inputenc}
\usepackage[T1]{fontenc}
\usepackage{lmodern}
\usepackage{amsmath,amssymb,amsthm,mathtools}
\usepackage[a4paper,margin=2.3cm]{geometry}
\usepackage{enumitem}
\usepackage{xcolor}
\usepackage[hidelinks]{hyperref}
\renewcommand{\leq}{\leqslant}
\renewcommand{\le}{\leqslant}
\renewcommand{\geq}{\geqslant}
\renewcommand{\ge}{\geqslant}
\setlength{\parindent}{0pt}
\setlength{\parskip}{0.35em}
% \hintstitle{Sheet title}{Class work / Homework}
\newcommand{\hintstitle}[2]{\begin{center}{\LARGE\bfseries Hints: #1\par}\vspace{0.5em}{\large #2\par}\end{center}
\noindent\textit{The hints to each problem are cumulative. Read one, try again, and open the next only if you are still stuck.}\par\medskip\hrule\medskip}
\newcommand{\problem}[1]{\par\bigskip{\large\bfseries Problem #1}\par\nopagebreak\smallskip}
\newcommand{\subpart}[1]{\par\smallskip\textbf{(#1)}\par\nopagebreak}
\newcommand{\hint}[1]{\par\smallskip\noindent\textcolor{black!70}{\textbf{Hint #1.}}\enspace}
\newcommand{\note}[1]{\par\smallskip\noindent\textit{Note.} #1\par}
\providecommand{\lt}{<}
\providecommand{\gt}{>}
\providecommand{\tightlist}{\setlength{\itemsep}{0pt}\setlength{\parskip}{0pt}}
\emergencystretch=3em

\begin{document}
\hintstitle{Combinatorics -- 2. Recurrences and generating functions}{Class work}

\problem{1}
\hint{1}
Multiply \(A(x)\) by \(Q(x)=1-c_1x-\dots-c_kx^k\): the coefficient of \(x^n\) in \(Q(x)A(x)\) is \(a_n-c_1a_{n-1}-\dots-c_ka_{n-k}\), which vanishes for \(n\ge k\). So \(QA\) is a polynomial of degree \(\lt k\), and conversely.
\hint{2}
For the explicit form, decompose \(P/Q\) into partial fractions \(\sum_i\sum_{j\lt m_i}\frac{b_{ij}}{(1-\lambda_ix)^{j+1}}\) and use \([x^n](1-\lambda x)^{-(j+1)}=\binom{n+j}{j}\lambda^n\), a polynomial in \(n\) of degree \(j\) times \(\lambda^n\). For \(a_n=4a_{n-1}-4a_{n-2}\) the root \(2\) is double, so \(a_n=(\alpha+\beta n)2^n\); the initial values give \(a_n=(n+2)2^{n-1}\).
\hint{3}
If all \(\lvert\lambda_i\rvert\lt1\), then \(a_n=\sum p_i(n)\lambda_i^n\to0\); integers tending to \(0\) are eventually \(0\). (Once \(k\) consecutive terms vanish, the recurrence keeps all later ones at \(0\) anyway.)

\problem{2}
\hint{1}
Solve \((\sum a_nx^n)(\sum b_nx^n)=1\) coefficient by coefficient: \(a_0b_0=1\) and \(a_0b_n=-\sum_{k\ge1}a_kb_{n-k}\) determine the \(b_n\) recursively exactly when \(a_0\ne0\).
\hint{2}
For \(1/(1-x-x^2)=\sum b_nx^n\) the recursion gives \(b_n=b_{n-1}+b_{n-2}\), \(b_0=b_1=1\), so \(b_n=F_{n+1}\). The dictionary rules: a pair of objects of sizes \(i,j\) contributes \(x^{i+j}\), so the generating function of pairs is the product; a sequence of \(m\) pieces is counted by \(F^m\), and \(\sum_m F^m=1/(1-F)\) is a well-defined formal series when \(F(0)=0\).
\hint{3}
Compositions into parts \(1\) and \(2\): sequences of pieces counted by \(F=x+x^2\), so the generating function is \(1/(1-x-x^2)\) and the count is \(F_{n+1}\).

\problem{3}
\hint{1}
Compute a few more terms or test a linear recurrence of order 2 on the given ones: \(26=6\cdot5-4\cdot1\), \(136=6\cdot26-4\cdot5\), \ldots{} The roots of \(t^2-6t+4\) are \(3\pm\sqrt5\).
\hint{2}
Define \(y_n=A(3+\sqrt5)^n+B(3-\sqrt5)^n\) with \(y_0=1\), \(y_1=5\), i.e.~\(A=\frac{5+2\sqrt5}{10}\), \(B=\frac{5-2\sqrt5}{10}\). Check that \(y_{n+1}-3y_n=\sqrt5\,(A\alpha^n-B\beta^n)\) while \(\sqrt5\,y_n=\sqrt5\,(A\alpha^n+B\beta^n)\).
\hint{3}
So \(\sqrt5\,y_n-(y_{n+1}-3y_n)=2\sqrt5\,B(3-\sqrt5)^n\in(0,1)\), hence \(y_{n+1}=3y_n+\lfloor y_n\sqrt5\rfloor\) and \(x_n=y_n=\frac{5+2\sqrt5}{10}(3+\sqrt5)^n+\frac{5-2\sqrt5}{10}(3-\sqrt5)^n\) by induction.

\problem{4}
\hint{1}
Look for a polynomial \(f\) of degree \(\le3\) with \(f(0)=1\) and \(f(k)\le0\) for all integers \(k\ge1\); then \(\Pr(X=0)\ge E f(X)\), and \(Ef(X)\) is computable from the moments.
\hint{2}
Take \(f(x)=-\frac{(x-1)(x-2)(x-3)}{6}\). Then \(Ef(X)=-\frac{EX^3-6EX^2+11EX-6}{6}=\frac13\). Equality needs \(X\) supported on \(\{0,1,2,3\}\): solve for the distribution (\(\Pr(X=0)=\frac13\), \(\Pr(X=1)=\frac12\), \(\Pr(X=3)=\frac16\)). The minimum is \(\frac13\).

\problem{5}
\hint{1}
Words without \(11\): states ``last letter \(0\)'' and ``last letter \(1\)'', matrix \(\begin{pmatrix}1&1\\1&0\end{pmatrix}\); the count is \(F_{n+2}\) and the growth rate \(\frac{1+\sqrt5}2\). For \(2\times n\) domino tilings, cut at the first column: either a vertical domino or two horizontal ones, giving \(t_n=t_{n-1}+t_{n-2}\).
\hint{2}
For \(3\times n\), use as states the set of cells of the next column already covered by horizontal dominoes sticking out of the current one (8 subsets, but only a few are reachable). Equivalently, \(t_m\) for a \(3\times2m\) board satisfies \(t_m=3t_{m-1}+2(t_{m-2}+\dots+t_0)\), i.e.~\(t_m=4t_{m-1}-t_{m-2}\), and odd widths give \(0\).
\hint{3}
So the generating function is \(\frac{1-x^2}{1-4x^2+x^4}\) in the width \(n\); characteristic polynomial \(t^2-4t+1\) in \(m\); growth rate \(2+\sqrt3\) per two columns, i.e.~\(\sqrt{2+\sqrt3}\approx1.932\) per column.
\end{document}
